\section{The projective comparison factor}
\label{compare:section}

We now extend the projective factor supplied by
Proposition~\ref{spread:product}.  The relevant positivity input is
algebraic b-semiampleness.  We apply it to a projective comparison
family, retaining the specified identification of its plurivolumes.

\begin{proposition}
\label{compare:semiample}
Use the notation of Proposition~\ref{spread:product}, and write
\(P^\circ\subset P\) for its good open set.  After a proper modification
\(\pi:P'\to P\), with \(P'\) smooth compact, and a further shrinking of
\(P^\circ\), there are a positive integer \(a\), a globally generated
holomorphic line bundle \(\mathcal H_Q\) on \(P'\), and a specified
isomorphism
\begin{equation}
\label{compare:open-identification}
 \mathcal H_Q|_{\pi^{-1}(P^\circ)}
 \simeq
 \bigl(\pi^*\mathcal G_{m,Q}\bigr)^{\otimes a}
       |_{\pi^{-1}(P^\circ)} .
\end{equation}
Transport to the left the \(a\)-th power of the volume norm
\[
 N_{m,Q}(s)=
 \left(\int_{Q_x}|s|^{2/m}\right)^{m/2}.
\]
For any holomorphic disk in \(P'\) whose punctured disk lies in the
good open, a local generator of the pulled-back \(\mathcal H_Q\)
and its dual have norm bounded by powers of
\(1+|\log|t||\), after a finite power substitution if necessary.
\end{proposition}

\begin{proof}
Let \(H\) be the product of the fixed Hilbert schemes used to record
the embedded projective factor and its boundary components.
Proposition~\ref{spread:product} gives a holomorphic map
\(\eta:P\to H\).
Its reduced image \(H_0\) is a compact irreducible analytic subspace.
Since \(H\) is projective, the proper mapping theorem and Chow's
theorem make \(H_0\) an irreducible projective variety.

We first identify an algebraic open of \(H_0\) carrying the desired
family.  Smoothness, connectedness of the fibers, the condition that
the marked subschemes are divisors, and their relative simple normal
crossing condition are algebraically constructible conditions on the
universal marked family.  The coefficients of the marks are the fixed
rational coefficients of \(C\), all in \([0,1)\).
On its smooth locus the bundle
\[
 H_m=\omega_{\mathrm{rel}}^{\otimes m}\otimes\cO(mC)
\]
is defined after the fixed choice of \(m\).  On a smooth connected
projective fiber, it is trivial precisely when
\[
 h^0(H_m)\geq1
 \quad\hbox{and}\quad
 h^0(H_m^{-1})\geq1.
\]
Indeed, the product of nonzero sections is a nonzero holomorphic
function and is therefore a nonzero constant.
These cohomology conditions are constructible by semicontinuity.
All the conditions hold on \(\eta(P^\circ)\), which is dense in
\(H_0\).  A constructible subset containing a dense subset of an
irreducible variety contains a nonempty algebraic open.
We may therefore choose an algebraic open \(H_0^\circ\) on which
the universal marked family is smooth projective, its fibers are
connected, and \(H_m\) is trivial on every fiber.
After shrinking the good open of \(P\), its projective factor is the
pullback of this family.
Cohomology and base change gives a line bundle
\(\mathcal V_m\) on \(H_0^\circ\).
Its evaluation map on the universal family is an isomorphism:
on every fiber it evaluates the one-dimensional space of sections
of a trivial line bundle.

Take a smooth projective model \(R\to H_0\), and resolve the
dominant component of the pulled-back universal family.  We obtain
a projective morphism
\[
 q:Z_Q\longrightarrow R
\]
with \(Z_Q\) smooth projective, unchanged over a smaller smooth
open \(R^\circ\subset R\).
The generic fiber is smooth and geometrically connected.
Consequently the finite map in the Stein factorization of \(q\)
is birational; it is an isomorphism because \(R\) is normal.
Thus \(q\) has connected fibers.

Choose a rational frame of \(\mathcal V_m\) over \(R\).
On \(Z_Q\) it defines a rational relative \(m\)-pluriform
\(\sigma\).  Shrink \(R^\circ\) so that this frame is regular and
nonvanishing there.  As a meromorphic section of
\(\omega_{Z_Q/R}^{\otimes m}\), its divisor has horizontal part
\(-mC\).  Define the rational divisor
\begin{equation}
\label{compare:volume-boundary}
 D_Q=-\frac1m\operatorname{div}(\sigma).
\end{equation}
The form \(\sigma\) then supplies an actual trivialization
\[
 \omega_{Z_Q/R}^{\otimes m}\otimes\cO(mD_Q)\simeq\cO_{Z_Q},
\]
and hence an adjoint identity
\begin{equation}
\label{compare:adjoint}
 K_{Z_Q}+D_Q\sim_{\Q}q^*K_R.
\end{equation}
Over \(R^\circ\), the divisor \(D_Q\) is exactly \(C\).
Its remaining coefficients are vertical and may have either sign.
Resolve their support, taking the boundary crepantly; this preserves
\eqref{compare:adjoint} and does not change the smooth family over
\(R^\circ\).

We check the algebraic b-semiampleness input explicitly.
Theorem~1.5 of \cite{BFMT25}, with Definition~6.18, applies to a
projective lc-trivial fibration whose boundary is effective over the
generic point.  The required conditions are generic sub-log-canonicity,
the discrepancy rank-one condition, and a rational adjoint pullback
identity.  Here the generic pair is the smooth effective klt pair
\((Q,C)\), and \eqref{compare:adjoint} is the required identity.
For the rank condition, on the generic fiber the rounded
discrepancy divisor is \(\lceil-C\rceil=0\); connectedness gives
\(h^0(Q,\cO_Q)=1\).  A resolution of the vertical locus does not
affect this generic computation.  Thus all these hypotheses hold.

For completeness, vertical adjustments do not change the moduli
object.  If a rational divisor \(E\) on \(R\) is added to the
adjoint base bundle and \(q^*E\) is added to \(D_Q\), then
\[
 t_P\longmapsto t_P-\operatorname{coeff}_P(E),
 \qquad
 B_R\longmapsto B_R+E,
 \qquad
 M_R\longmapsto M_R.
\]
The same equalities hold on every higher base model.
Thus one may, if desired, arrange upper bounds on the vertical
coefficients by subtracting sufficiently large pullbacks.
Alternatively one may arrange global effectivity by adding such
pullbacks.  Each operation is possible because the finitely many
vertical components have proper algebraic images contained in
divisors on the projective base.  Neither operation is needed:
the theorem imposes effectivity and log canonicity over the generic
point.

It follows from Theorem~1.5 of \cite{BFMT25} that the moduli
b-divisor is b-semiample.  Its identification with the moduli object
defined by the log canonical thresholds is supplied by
Theorem~6.28 of \cite{BFMT25}: the Hodge-theoretic and threshold
moduli divisors agree up to a fixed rational linear equivalence.
Consequently, on a sufficiently high smooth projective model
\(\mu_R:R'\to R\), the actual line bundle
\[
 \mathcal L_Q=\cO_{R'}(amM_{R'})
\]
is globally generated for some positive integer \(a\).
Increase \(a\) to clear all denominators.
Further projective modifications preserve global generation of the
pullback, so we may also arrange a simple normal crossing complement
to the good open.

We must retain the particular open identification with plurivolumes.
On a smooth source model over \(R'\), take the crepant transform
\(D'_Q\) and the base bundle \(L_{R'}=\mu_R^*K_R\).
Then
\[
 K_{Z'_Q/R'}+D'_Q
 \sim_{\Q}
 (q')^*(L_{R'}-K_{R'}).
\]
Over the common good open this identity identifies the direct image
of the log \(m\)-pluricanonical bundle with the line of the original
volumes.  Theorem~\ref{orders:trace} applies to this projective
family: its generic boundary is effective klt, and the theorem
permits arbitrary vertical coefficients.
It identifies the extension specified by the threshold trace
\(mM_{R'}\) under precisely this open-volume identification.
The base Jacobian on \(R'\) is included in that theorem's relative
order calculation.  In particular, the \(k=0\) case of
Corollary~\ref{orders:residuenorm} gives the two-sided logarithmic
norm bounds for local generators and their duals, including after
pullback to a disk.

Finally resolve the graph of the meromorphic lift
\(P\dashrightarrow R'\) of \(\eta\), together with the complement
of the good open, obtaining \(\pi:P'\to P\) and
\(\rho:P'\to R'\), with \(P'\) smooth compact and its boundary
simple normal crossing.
Set
\[
 \mathcal H_Q=\rho^*\mathcal L_Q.
\]
This line bundle is globally generated.  The preceding specified
identification, pulled back along \(\rho\), is
\eqref{compare:open-identification}, and the same pullback gives the
asserted norm bounds.
\end{proof}

\begin{remark}
The rational form in \eqref{compare:volume-boundary} is taken on the
projective comparison space \(Z_Q\).  The construction makes no
assumption that the canonical bundle, or an arbitrary holomorphic
line bundle, on the original nonprojective manifold has a global
meromorphic section.  It also specifies an isomorphism of actual
bundles over the open set; a numerical identification would not
suffice for the final comparison.
\end{remark}
